% generated by tools/make_tex_tables.py -- do not edit
% The paragraph that belongs to tables/table9_surveys.tex: what the columns mean and how they were
% read. It is body prose, input by the section that owns the table, and not a second
% block of small type inside the float.
A cell is a level and not a score: \cvyes\ where the work gives the topic a structure of its own, a section, a table or a measurement it defines; \cvpart\ where the topic is there without one, a subsection with no analysis, a handful of mentions, or a metric named and never defined; \cvno\ where it is absent or put outside the work's scope. Each level coarsens one sentence read from \path{papers/notes/<key>.md}, and it is printed only while the quotation behind that sentence is still in the note, so a glyph cannot outlive its evidence. The supporting sentences remain in the reproducibility materials and \path{paper/tables/table10_surveys.md}. Three of the fourteen sources could not be obtained and one is on disk and has never been read, which is why their cells are empty; they are grouped as such rather than scored, because this survey knows nothing of them beyond a title and a DOI. The \emph{gaps named} column is uniform by design: every predecessor names gaps, and what none of them attaches to a gap is a measurement.
